%% -----------------------------------------------------------------
%% 北京邮电大学研究生学位论文 - 英文摘要
%% 格式要求（2025年1月15日规范）：
%% 1. “ABSTRACT”标题用 Times New Roman 字体三号字加粗，居中对齐，首行无缩进，段前0行，段后2行，单倍行距
%% 2. 摘要内容用四号 Times New Roman 字体，两端对齐，首行缩进2字符，段前0行，段后0行，固定值20磅
%% 3. “KEY WORDS:”四号字加粗，其后关键词为四号字，两端对齐，首行缩进2字符，段前2行，段后0行，固定值20磅
%% 4. 英文关键词与中文关键词对应，每一关键词之间用半角分号隔开（; ），最后一个关键词后不打标点符号
%% -----------------------------------------------------------------

\begin{enabstract}
The abstract should concisely present the essential contents and core methodology of the thesis. It highlights the research background, objectives, methodologies, significant theoretical contributions, experimental results, and conclusions.

This thesis focuses on key scientific challenges and engineering problems in next-generation information and communication networks. First, the existing literature and fundamental bottlenecks in the field are comprehensively summarized. Second, a novel adaptive optimization framework and distributed processing mechanism are proposed to achieve high robustness and low processing delay. Third, rigorous theoretical derivations are provided to characterize the convergence speed and capacity bounds of the proposed algorithms. Finally, extensive simulation studies and experimental evaluations on real-world datasets demonstrate the superior efficiency and reliability of the proposed scheme.

The results show that the proposed approach substantially reduces computational complexity while improving throughput and system accuracy, validating its effectiveness in realistic complex environments.

\enkeywords{Beijing University of Posts and Telecommunications; Graduate Degree Thesis; \LaTeX{} Template; Formatting Specification; Information and Communication Engineering}
\end{enabstract}
